%% ============================================================
%% Diapositivas Práctica 04 — Dog Messenger
%% CC4P1 Programación Concurrente y Distribuida  2026-I
%% Escuela de Ciencias de la Computación — UNI
%% ============================================================
\documentclass[aspectratio=169,14pt]{beamer}

\usetheme{Madrid}
\usecolortheme{default}

\usepackage[spanish,es-noshorthands]{babel}
\usepackage[utf8]{inputenc}
\usepackage[T1]{fontenc}
\usepackage{listings}
\usepackage{xcolor}
\usepackage{graphicx}
\usepackage{booktabs}
\usepackage{hyperref}
\usepackage{tikz}
\usepackage{amssymb}
\usepackage{amsmath}
\usepackage{makecell}
\usetikzlibrary{shapes.geometric, arrows.meta, positioning}

%% ── Colores ──────────────────────────────────────────────────────────────────
\definecolor{accentblue}{RGB}{30,90,200}
\definecolor{accentgreen}{RGB}{20,140,60}
\definecolor{accentorange}{RGB}{200,100,0}
\definecolor{javabg}{RGB}{248,248,255}
\definecolor{kotlinbg}{RGB}{248,255,248}
\definecolor{pybg}{RGB}{255,252,240}
\definecolor{commentcolor}{RGB}{100,120,100}
\definecolor{keywordcolor}{RGB}{0,0,180}
\definecolor{stringcolor}{RGB}{160,32,32}
\definecolor{framecolor}{RGB}{180,180,200}

\setbeamercolor{title}{fg=white,bg=accentblue}
\setbeamercolor{frametitle}{fg=white,bg=accentblue}

%% ── Estilos listings ─────────────────────────────────────────────────────────
\lstdefinestyle{java}{
  language=Java,
  backgroundcolor=\color{javabg},
  basicstyle=\ttfamily\scriptsize,
  keywordstyle=\color{keywordcolor}\bfseries,
  commentstyle=\color{commentcolor}\itshape,
  stringstyle=\color{stringcolor},
  breaklines=true, breakatwhitespace=true,
  frame=single, rulecolor=\color{framecolor},
  tabsize=2, showstringspaces=false,
  morekeywords={var,val,record,sealed,permits,yield}
}

\lstdefinestyle{kotlin}{
  language=Java,
  backgroundcolor=\color{kotlinbg},
  basicstyle=\ttfamily\scriptsize,
  keywordstyle=\color{accentgreen}\bfseries,
  commentstyle=\color{commentcolor}\itshape,
  stringstyle=\color{stringcolor},
  breaklines=true, breakatwhitespace=true,
  frame=single, rulecolor=\color{framecolor},
  tabsize=2, showstringspaces=false,
  morekeywords={fun,val,var,data,class,object,companion,override,suspend,launch,coroutineScope}
}

\lstdefinestyle{python}{
  language=Python,
  backgroundcolor=\color{pybg},
  basicstyle=\ttfamily\scriptsize,
  keywordstyle=\color{accentorange}\bfseries,
  commentstyle=\color{commentcolor}\itshape,
  stringstyle=\color{stringcolor},
  breaklines=true, breakatwhitespace=true,
  frame=single, rulecolor=\color{framecolor},
  tabsize=2, showstringspaces=false
}

\lstset{
  literate=
    {á}{{\'{a}}}1 {é}{{\'{e}}}1 {í}{{\'{i}}}1 {ó}{{\'{o}}}1 {ú}{{\'{u}}}1
    {Á}{{\'{A}}}1 {É}{{\'{E}}}1 {Í}{{\'{I}}}1 {Ó}{{\'{O}}}1 {Ú}{{\'{U}}}1
    {ñ}{{\~{n}}}1 {Ñ}{{\~{N}}}1 {ü}{{\"u}}1    {Ü}{{\"U}}1
}

%% ── Comandos auxiliares ──────────────────────────────────────────────────────
\newcommand{\captura}[2]{
  \begin{figure}
    \centering
    \fcolorbox{gray}{white}{\begin{minipage}{0.85\textwidth}
      \centering
      \vspace{3cm}
      {\Large \color{gray!60} [Captura: #2]}
      \vspace{3cm}
    \end{minipage}}
    \caption{#2}
  \end{figure}
}

\newcommand{\nododiagram}[2]{
  \begin{figure}
    \centering
    \fcolorbox{gray}{white}{\begin{minipage}{0.7\textwidth}
      \centering
      \vspace{2.5cm}
      {\Large \color{gray!60} [Diagrama: #2]}
      \vspace{2.5cm}
    \end{minipage}}
    \caption{#2}
  \end{figure}
}

%% =============================================================================
\begin{document}

%% ── Portada ──────────────────────────────────────────────────────────────────
\begin{frame}
\titlepage
\end{frame}

%% ── Índice ───────────────────────────────────────────────────────────────────
\begin{frame}{Contenido}
\tableofcontents
\end{frame}

%% =============================================================================
\section{Introducción}

\begin{frame}{Dog Messenger}
\begin{itemize}
  \item Sistema de mensajería distribuida en tiempo real
  \item Clientes heterogéneos: Java desktop, Android/Kotlin, Python bot
  \item Comunicación via sockets TCP binarios (sin frameworks)
  \item Concurrencia mediante hilos, estructuras thread-safe
  \item Encriptación punto a punto (ECDH + AES-256)
  \item Nodo de ventas automatizado con NLP
\end{itemize}

\vfill
\centering
\emph{Tres lenguajes, un solo protocolo}
\end{frame}

%% =============================================================================
\section{Arquitectura del Sistema}

\begin{frame}{Arquitectura Cliente-Servidor (Estrella)}
\centering
\begin{tikzpicture}[
  node distance=1.5cm and 2cm,
  box/.style={rectangle, rounded corners=4pt, draw, minimum width=2.4cm,
              minimum height=0.9cm, align=center, font=\small},
  server/.style={box, fill=blue!15, draw=accentblue, thick},
  client/.style={box, fill=green!10, draw=accentgreen},
  bot/.style={box, fill=orange!15, draw=accentorange},
  arr/.style={-Latex, thick}
]
\node[server] (srv) {\textbf{Servidor}\\Java 17};

\node[client, above left=of srv]  (dc1) {Desktop\\Java};
\node[client, above right=of srv] (dc2) {Desktop\\Java};
\node[client, below left=of srv]  (and) {Android\\Kotlin};
\node[bot,    below right=of srv] (bot) {Nodo Ventas\\Python};

\draw[arr,<->] (srv) -- (dc1) node[midway,above left,font=\tiny]{TCP};
\draw[arr,<->] (srv) -- (dc2) node[midway,above right,font=\tiny]{TCP};
\draw[arr,<->] (srv) -- (and) node[midway,below left,font=\tiny]{TCP};
\draw[arr,<->] (srv) -- (bot) node[midway,below right,font=\tiny]{TCP};
\end{tikzpicture}
\end{frame}

\begin{frame}{Patrones de Concurrencia}
\begin{description}
  \item[Thread-per-Connection] \hfill \\
    Cada cliente ejecuta su propio \texttt{ClientHandler} en un hilo separado.
  \item[Producer-Consumer] \hfill \\
    \texttt{BlockingQueue} separa la GUI del envío al socket.
  \item[Thread-safe Registry] \hfill \\
    \texttt{ConcurrentHashMap} + \texttt{AtomicInteger} para registro concurrente.
\end{description}

\vfill
\nododiagram{width=0.6\textwidth}{Flujo de hilos en el servidor}
\end{frame}

%% =============================================================================
\section{Protocolo Binario}

\begin{frame}{Protocolo Binario --- Header de 9 bytes}
\begin{center}
\small
\begin{tabular}{|c|c|c|c|}
\hline
\textbf{payload\_len} & \textbf{type} & \textbf{sender\_id} & \textbf{receiver\_id} \\
4 bytes (int) & 1 byte & 2 bytes (short) & 2 bytes (short) \\
\hline
\end{tabular}
\end{center}

\vfill
\begin{itemize}
  \item Big-endian, compatible con \texttt{DataInputStream/OutputStream}
  \item 11 tipos de mensaje (AUTH, TEXT, FILE, GROUP, QR, SALES, METRICS, DH\_EXCHANGE, ...)
  \item Misma implementación en Java, Kotlin y Python
\end{itemize}

\vfill
\centering
\emph{Sin JSON, sin XML, sin WebSocket --- solo bytes crudos}
\end{frame}

\begin{frame}[fragile]{Parser en los tres lenguajes}
\begin{columns}[T]
\column{0.33\textwidth}
\centering{\textbf{Java}}
\lstinputlisting[style=java,firstline=326,lastline=334]{informe.tex}
\column{0.33\textwidth}
\centering{\textbf{Kotlin}}
\lstinputlisting[style=kotlin,firstline=617,lastline=625]{informe.tex}
\column{0.33\textwidth}
\centering{\textbf{Python}}
\begin{lstlisting}[style=python]
import struct

def read_header(sock):
    data = sock.recv(9, MSG_WAITALL)
    payload_len, type_code, \
    sender_id, receiver_id = \
        struct.unpack('>iBHH', data)
    return (type_code,
            sender_id, receiver_id,
            payload_len)
\end{lstlisting}
\end{columns}
\end{frame}

%% =============================================================================
\section{Servidor}

\begin{frame}[fragile]{Servidor --- Thread-per-Connection}
\begin{lstlisting}[style=java]
try (ServerSocket server = new ServerSocket(port)) {
    while (true) {
        Socket client = server.accept();
        new ClientHandler(client).start();
    }
}
\end{lstlisting}

\vfill
\begin{itemize}
  \item Hilo principal solo acepta conexiones
  \item Cada \texttt{ClientHandler} corre en su propio hilo
  \item Dispatch por tipo de mensaje via \texttt{switch}
  \item Envío \texttt{synchronized} para evitar entrelazado de bytes
\end{itemize}
\end{frame}

\begin{frame}[fragile]{Registro Thread-safe}
\begin{lstlisting}[style=java]
public class UserRegistry {
    private final ConcurrentHashMap<String, Short> usernames;
    private final AtomicInteger idCounter = new AtomicInteger(1);

    public short register(String username, ClientHandler handler) {
        short newId = (short) idCounter.getAndIncrement();
        if (usernames.putIfAbsent(username, newId) != null)
            return -1;
        handlers.put(newId, handler);
        return newId;
    }
}
\end{lstlisting}

\begin{itemize}
  \item \texttt{putIfAbsent} atómico --- evita duplicados sin \texttt{synchronized}
  \item \texttt{AtomicInteger} --- IDs únicos sin condición de carrera
\end{itemize}
\end{frame}

%% =============================================================================
\section{Cliente Desktop}

\begin{frame}[fragile]{Patrón Producer-Consumer}
\begin{lstlisting}[style=java]
public class SenderThread extends Thread {
    private final BlockingQueue<QueuedMessage> queue =
                                    new LinkedBlockingQueue<>();

    public void queueMessage(Message header, byte[] payload) {
        queue.offer(new QueuedMessage(header, payload));
    }

    public void run() {
        while (!shutdown) {
            QueuedMessage msg = queue.take();
            if (crypto != null && crypto.hasKeyFor(...))
                data = crypto.encrypt(receiverId, data);
            MessageParser.writeMessage(out, msg.header(), data);
        }
    }
}
\end{lstlisting}

\begin{itemize}
  \item GUI nunca bloquea: solo encola
  \item Encriptación E2E transparente antes de enviar
\end{itemize}
\end{frame}

\begin{frame}[fragile]{Encriptación Punto a Punto}
\begin{columns}
\column{0.55\textwidth}
\begin{itemize}
  \item Intercambio ECDH (secp256r1) via mensaje \texttt{DH\_EXCHANGE}
  \item Secreto compartido $\rightarrow$ SHA-256 $\rightarrow$ clave AES-256
  \item Cifrado AES-256/CBC con IV aleatorio
  \item El servidor solo ve bytes cifrados
\end{itemize}
\column{0.45\textwidth}
\centering
\small
\includegraphics[width=\textwidth]{placeholder-e2e.png}
\end{columns}
\end{frame}

\begin{frame}{Capturas --- Cliente Desktop}
\captura{width=0.9\textwidth}{Ventana principal de chat con mensajes de texto e imagen inline}

\captura{width=0.9\textwidth}{Conversación con el nodo de ventas (flujo de pedido completado)}
\end{frame}

%% =============================================================================
\section{Cliente Móvil Android}

\begin{frame}[fragile]{Android --- Mismo Protocolo, Mismo Código}
\begin{lstlisting}[style=kotlin]
object MessageParser {
    fun readHeader(input: DataInputStream): Message {
        val payloadLen  = input.readInt()
        val typeCode    = input.readByte().toInt() and 0xFF
        val senderId    = input.readShort()
        val receiverId  = input.readShort()
        return Message(payloadLen, ...)
    }
}
\end{lstlisting}

\begin{itemize}
  \item \texttt{DataInputStream/OutputStream} de la JVM
  \item Compatibilidad byte a byte con el servidor Java
  \item Envío protegido con \texttt{synchronized(this)}
  \item Guardado de archivos en \texttt{MediaStore} (API 30+)
\end{itemize}
\end{frame}

\begin{frame}{Capturas --- Cliente Android}
\captura{width=0.5\textwidth}{Pantalla de chat en dispositivo Android}
\end{frame}

%% =============================================================================
\section{Nodo de Ventas}

\begin{frame}[fragile]{Nodo de Ventas (Python)}
\begin{itemize}
  \item Se conecta como cliente normal, username \texttt{"ventas"}
  \item El servidor le reenvía mensajes SALES de forma transparente
  \item Motor NLP con similitud coseno (bag-of-words, sin librerías externas)
  \item Máquina de estados conversacional por usuario
\end{itemize}

\begin{lstlisting}[style=python]
class SalesNode:
    def _authenticate(self):
        send_message(self.sock, AUTH, 0, BROADCAST, b'ventas')

    def _dispatch(self, type_code, sender_id, payload):
        if type_code == TEXT:
            response = self.chatbot.handle(sender_id, text)
\end{lstlisting}
\end{frame}

\begin{frame}[fragile]{Motor NLP --- Similitud Coseno}
\begin{columns}[T]
\column{0.55\textwidth}
\begin{lstlisting}[style=python]
import math, re
from collections import Counter

def cosine_sim(v1, v2):
    common = set(v1) & set(v2)
    dot = sum(v1[k] * v2[k] for k in common)
    mag1 = math.sqrt(sum(v*v for v in v1.values()))
    mag2 = math.sqrt(sum(v*v for v in v2.values()))
    return dot / (mag1 * mag2) if mag1 and mag2 else 0.0
\end{lstlisting}
\column{0.45\textwidth}
\small
\textbf{Corpus:}
\begin{itemize}
  \item 802 stopwords
  \item 324 términos de productos
  \item 172 frases en 5 intenciones
  \item 667 tokens significativos
\end{itemize}
\end{columns}
\end{frame}

\begin{frame}{Ejemplo de Conversación}
\centering
\small
\begin{tabular}{p{4.5cm} p{6cm}}
\toprule
\textbf{Usuario} & \textbf{Bot (intención)} \\
\midrule
cuales son mis pedidos &
  CONSULTAR (score=0.35) --- \textit{Cuál es tu ID?} \\[4pt]
quiero registrarme &
  REGISTRAR (score=0.62) --- \textit{Cuál es tu nombre?} \\[4pt]
quisiera comprar croquetas &
  PEDIDO (score=0.50) --- \textit{Numero de cliente?} \\[4pt]
dame el reporte de ventas &
  REPORTE (score=0.53) --- \textit{Estadísticas} \\
\bottomrule
\end{tabular}
\end{frame}

\begin{frame}{Captura --- Nodo de Ventas}
\captura{width=0.85\textwidth}{Consola del nodo de ventas mostrando autenticación y respuestas NLP}
\end{frame}

%% =============================================================================
\section{Métricas de Rendimiento}

\begin{frame}[fragile]{Métricas Lock-Free}
\begin{lstlisting}[style=java]
public class MetricsCollector {
    private final AtomicInteger totalConnections = new AtomicInteger(0);
    private final AtomicInteger totalMessages    = new AtomicInteger(0);
    private final AtomicLong    totalBytes       = new AtomicLong(0);

    public void recordMessage(int payloadLen) {
        totalMessages.incrementAndGet();
        totalBytes.addAndGet(payloadLen);
    }
}
\end{lstlisting}

\begin{itemize}
  \item \texttt{Atomic*} sin bloqueos --- operaciones CAS del procesador
  \item Sin contención bajo alta concurrencia
  \item Cliente obtiene métricas con mensaje tipo METRICS
\end{itemize}
\end{frame}

%% =============================================================================
\section{Clonación por QR}

\begin{frame}{Clonación de Chat por QR}
\centering
\begin{enumerate}
  \item Usuario A solicita token QR desde el desktop
  \item El servidor genera UUID con TTL de 120s
  \item Se renderiza como QR (ZXing)
  \item Usuario B escanea el código
  \item Servidor valida y envía el historial a B
\end{enumerate}

\vfill
\nododiagram{width=0.7\textwidth}{Flujo de clonación QR entre dispositivos}
\end{frame}

%% =============================================================================
\section{Conclusiones}

\begin{frame}{Conclusiones}
\begin{itemize}
  \item Sistema completo usando solo sockets crudos e hilos
  \item Protocolo binario de 9 bytes --- clave de interoperabilidad
  \item \texttt{ConcurrentHashMap}, \texttt{AtomicInteger}, \texttt{BlockingQueue}
        eliminan bloqueos globales
  \item Nodo Python se integra sin modificar el servidor
  \item Encriptación E2E sin depender del servidor
\end{itemize}

\vfill
\centering
\emph{Un protocolo, tres lenguajes, cero dependencias externas de comunicación}
\end{frame}

\end{document}
